% BHU PYQ -- Transcribed & Corrected
% Paper: MATMJ-44 : Mechanics
% Exam: B.Sc. (Hons.) Semester IV Examination, 2025-26 (NEP)
% Department: Department of Mathematics

\documentclass[11pt,a4paper]{article}
\usepackage[a4paper,top=2.5cm,bottom=2.5cm,left=2.8cm,right=2.8cm,headheight=14pt]{geometry}
\usepackage{amsmath,amssymb}
\usepackage[T1]{fontenc}
\usepackage[utf8]{inputenc}
\usepackage{lmodern,microtype,fancyhdr,enumitem,hyperref,lastpage,parskip}

\hypersetup{
    colorlinks=true,
    urlcolor=black,
    linkcolor=black,
    pdftitle={MATMJ44: Mechanics -- BHU B.Sc. Sem IV 2025-26 (NEP)}
}

\pagestyle{fancy}
\fancyhf{}
\renewcommand{\headrulewidth}{0.4pt}
\renewcommand{\footrulewidth}{0.4pt}
\fancyhead[L]{\small\textit{Banaras Hindu University}}
\fancyhead[R]{\small\textit{MATMJ-44: Mechanics}}
\fancyfoot[C]{\small Page \thepage\ of \pageref{LastPage}}

\newlist{parts}{enumerate}{2}
\setlist[parts,1]{label=(\roman*),leftmargin=*,itemsep=4pt,topsep=2pt}
\newcommand{\pts}[1]{\hfill\textit{[#1]}}

\begin{document}

\begin{center}
  {\large\bfseries BANARAS HINDU UNIVERSITY}\\[2pt]
  {\normalsize B.Sc. (Hons.) Semester IV Examination, 2025-26 (NEP)}\\[6pt]
  \rule{0.8\linewidth}{0.5pt}\\[6pt]
  {\large\bfseries DEPARTMENT OF MATHEMATICS}\\[4pt]
  {\large\bfseries Paper Code: MATMJ-44 : Mechanics}\\[4pt]
  {\normalsize Time: 3 Hours\hfill Maximum Marks: 60}
\end{center}

\rule{\linewidth}{0.4pt}

\smallskip
\noindent\textit{Instructions: Attempt FIVE questions in total. All questions carry equal marks (12 marks each).}

\bigskip

\noindent\textbf{Question 1.} State and prove the Principle of Virtual Work for a system of coplanar forces acting on a rigid body.\pts{12}

\medskip\rule{\linewidth}{0.2pt}\medskip

\noindent\textbf{Question 2.} Derive the intrinsic equation of a Common Catenary $s = c \tan \psi$, and obtain relation $y = c \cosh(x/c)$.\pts{12}

\medskip\rule{\linewidth}{0.2pt}\medskip

\noindent\textbf{Question 3.} Derive the differential equation of a central orbit in pedal form $\frac{h^2}{p^3} \frac{dp}{dr} = F(r)$.\pts{12}

\medskip\rule{\linewidth}{0.2pt}\medskip

\noindent\textbf{Question 4.} State Kepler's laws of planetary motion and deduce Newton's Law of Gravitation from Kepler's laws.\pts{12}

\medskip\rule{\linewidth}{0.2pt}\medskip

\noindent\textbf{Question 5.} Derive tangential and normal accelerations of a particle moving along a plane curve.\pts{12}

\vfill
\rule{\linewidth}{0.4pt}
\begin{center}\small
\href{https://bhu-exam.vercel.app/}{Click here to visit our website}\\[4pt]
\href{https://me-aryan.vercel.app/}{Click here to visit our contributor}
\end{center}

\end{document}
